% Chapter 1: Abstract
% ============================================================
% (This file is included via \input in the abstract environment of main.tex)

We introduce and systematically study the \emph{Vertex-Distance Perimeter Problem}: given a triangle $\tri ABC$ in the Euclidean plane and an arbitrary interior point $P$, with known Euclidean distances $d_1 = PA$, $d_2 = PB$, and $d_3 = PC$ to the three vertices, determine the perimeter $\perim = a + b + c$ of the triangle, where $a$, $b$, $c$ denote the side lengths opposite vertices $A$, $B$, $C$ respectively.

We establish through rigorous degrees-of-freedom analysis that the problem, as stated, is \emph{fundamentally underdetermined}: a triangle possesses three intrinsic shape parameters (its side lengths), while an interior point contributes two additional positional degrees of freedom, yielding a system with five unknowns constrained by only three measurements. The resulting solution space forms a \emph{one-dimensional manifold} of geometrically distinct triangles after quotienting by the Euclidean group, parameterizable by a single angular parameter at the interior point.

The central mathematical apparatus is the \emph{Cayley--Menger determinant}, a $(n+2) \times (n+2)$ bordered symmetric matrix whose vanishing characterizes the realizability of $n+1$ points in $\R^{n-1}$. Applied to our four coplanar points $(A, B, C, P)$, the $5 \times 5$ Cayley--Menger determinant yields a polynomial constraint linking all six pairwise distances. We derive this constraint explicitly and analyze the resulting algebraic variety.

For constrained triangle classes, the problem admits finite or unique solutions. We derive closed-form solutions for equilateral triangles (via the identity $3(x^4 + y^4 + z^4 + s^4) = (x^2 + y^2 + z^2 + s^2)^2$), specialized formulas for isosceles and right triangles, and classification criteria for acute and obtuse configurations. We prove Pompeiu's theorem as a structural result connecting our framework to classical triangle geometry, and establish connections to Stewart's theorem and the Leibniz distance formula.

We present five distinct computational approaches---symbolic algebraic solving, Newton--Raphson iteration, gradient descent optimization, the Nelder--Mead simplex method, and hybrid bisection--secant root-finding---implemented across four programming languages (Python, TypeScript, Go, and Rust/WebAssembly). A synthetic dataset of $10^6$ solved triangle configurations supports extensive benchmarking of solver accuracy, convergence rates, and computational performance.

The theory is extended to non-Euclidean geometries: spherical triangles on $S^2$ via the spherical law of cosines and great-circle distances, and hyperbolic triangles in the Poincar\'e disk model via the hyperbolic metric. A three-dimensional generalization to tetrahedra is formulated using the $6 \times 6$ Cayley--Menger determinant.

An interactive web application provides real-time geometric exploration, solver visualization, error landscape rendering, and step-by-step proof walkthroughs.

\medskip
\noindent\textbf{Keywords:} Triangle reconstruction, vertex distances, interior point, Cayley--Menger determinant, Pompeiu's theorem, distance geometry, solution manifold, computational geometry, non-Euclidean geometry.

\noindent\textbf{MSC 2020:} 51M04, 52B11, 51K05, 65H10, 68U05.
